% Configuración compartida por KalaScript.tex y atlas.tex.
\usepackage[T1]{fontenc}
\usepackage[utf8]{inputenc}
\usepackage{lmodern}
\usepackage{babel}
\babelprovide[main,import]{spanish}
\usepackage[a4paper,margin=22mm,headheight=16pt]{geometry}
\usepackage{microtype}
\usepackage{amsmath,amssymb}
\usepackage{graphicx,xcolor,tikz}
\usepackage{booktabs,longtable,array,tabularx}
\usepackage{enumitem,fancyvrb,fancyhdr,titlesec}
\usepackage{xurl}
\usepackage[unicode,hidelinks]{hyperref}

\definecolor{KalaTinta}{HTML}{263B35}
\definecolor{KalaSuave}{HTML}{616C66}
\definecolor{KalaAcento}{HTML}{A75735}
\definecolor{KalaLinea}{HTML}{D7DDD3}
\definecolor{KalaPapel}{HTML}{F7F5EE}
\definecolor{KalaVerde}{HTML}{E6EBE1}
\hypersetup{pdftitle={KalaScript · Núcleo escrito},
  pdfsubject={Escritura visual, gramática y ejemplos},
  colorlinks=true,linkcolor=KalaTinta,urlcolor=KalaAcento,
  bookmarksnumbered=true}
\setlength{\parindent}{0pt}
\setlength{\parskip}{6pt plus 1pt}
\setlength{\emergencystretch}{2em}
\setlist{itemsep=3pt,topsep=5pt,leftmargin=*}
\renewcommand{\arraystretch}{1.15}
\renewcommand{\contentsname}{Contenido}
\renewcommand{\tablename}{Tabla}
\renewcommand{\figurename}{Figura}
\titleformat{\section}{\Large\sffamily\color{KalaTinta}}{\thesection}{.7em}{}
\titleformat{\subsection}{\large\sffamily\color{KalaTinta}}{\thesubsection}{.7em}{}
\titleformat{\subsubsection}{\normalsize\sffamily\bfseries\color{KalaTinta}}{\thesubsubsection}{.7em}{}
\pagestyle{fancy}
\fancyhf{}
\fancyhead[L]{\small\sffamily\color{KalaSuave}KalaScript}
\fancyhead[R]{\small\sffamily\color{KalaSuave}Núcleo escrito}
\fancyfoot[L]{\footnotesize\sffamily\color{KalaSuave}Escritura personal}
\fancyfoot[R]{\small\sffamily\thepage}
\renewcommand{\headrulewidth}{.3pt}
\renewcommand{\footrulewidth}{0pt}

% Registro único: código, nombre, familia, definición, trayectos, dibujo.
\newcommand{\KalaRegistrar}[6]{%
  \expandafter\def\csname KalaNombre#1\endcsname{#2}%
  \expandafter\def\csname KalaFamilia#1\endcsname{#3}%
  \expandafter\def\csname KalaDescripcion#1\endcsname{#4}%
  \expandafter\def\csname KalaTrayectos#1\endcsname{#5}%
  \expandafter\def\csname KalaDibujo#1\endcsname{#6}%
}
\newcommand{\KalaSimbolo}[2][9mm]{%
  \begingroup
  \ifcsname KalaDibujo#2\endcsname\else
    \PackageError{KalaScript}{Simbolo desconocido: #2}{Use un codigo del catalogo.}%
  \fi
  \pgfmathsetlengthmacro{\KalaUnidad}{#1/64}%
  \pgfmathsetlengthmacro{\KalaGrosor}{#1*2.6/64}%
  \begin{tikzpicture}[x=\KalaUnidad,y=-\KalaUnidad,
    baseline={([yshift=-.55ex]current bounding box.center)},
    draw=KalaTinta,line width=\KalaGrosor,line cap=round,line join=round]
    \path[use as bounding box] (0,0) rectangle (64,64);
    \csname KalaDibujo#2\endcsname
  \end{tikzpicture}%
  \endgroup
}
\newcommand{\KalaGrupo}[1]{%
  \tikz[baseline={([yshift=-.55ex]current bounding box.center)}]
    \node[draw=KalaSuave,line width=.35pt,rounded corners=.6pt,
      inner xsep=4pt,inner ysep=4pt]{#1};%
}
\newcommand{\KalaMat}[1]{%
  \tikz[baseline={([yshift=-.55ex]current bounding box.center)}]
    \node[draw=KalaSuave,dashed,line width=.35pt,fill=KalaPapel,
      inner xsep=4pt,inner ysep=3pt]{\ensuremath{\displaystyle #1}};%
}
\newcommand{\KalaReferencia}[1]{\KalaSimbolo{K10}\kern-.8mm\textsubscript{#1}}
\newcommand{\KalaSeparador}{\kern1.4mm\textcolor{KalaSuave}{\ensuremath{\mid}}\kern1.4mm}
\newcommand{\KalaRelacion}[2]{\KalaSimbolo{K03}\kern-.8mm\textsubscript{#1}\kern1.5mm\KalaGrupo{#2}}
\newcommand{\KalaNegacion}[1]{\KalaSimbolo{K11}\kern1mm\KalaGrupo{#1}}
\newcommand{\KalaConjuncion}[2]{\KalaSimbolo{K12}\kern1mm\KalaGrupo{#1\KalaSeparador#2}}
\newcommand{\KalaCondicional}[2]{\KalaSimbolo{K13}\kern1mm\KalaGrupo{#1\KalaSeparador#2}}
\newcommand{\KalaUniversal}[3]{\KalaSimbolo{K14}\kern1mm\ensuremath{#1\in}\,\KalaMat{#2}\kern1mm\KalaGrupo{#3}}
\newcommand{\KalaExistencia}[3]{\KalaSimbolo{K15}\kern1mm\ensuremath{#1\in}\,\KalaMat{#2}\kern1mm\KalaGrupo{#3}}
\newcommand{\KalaFuente}[2]{\KalaSimbolo{K16}\kern1mm#1\kern1mm\KalaGrupo{#2}}
\newcommand{\KalaDefinicion}[2]{\KalaSimbolo{K09}\kern1mm\ensuremath{#1}\kern1mm\KalaGrupo{#2}}
\newsavebox{\KalaCajaExpresion}
\newcommand{\KalaMostrar}[1]{%
  \begin{center}
  \sbox{\KalaCajaExpresion}{#1}%
  \ifdim\wd\KalaCajaExpresion>\linewidth
    \resizebox{\linewidth}{!}{\usebox{\KalaCajaExpresion}}%
  \else\usebox{\KalaCajaExpresion}\fi
  \end{center}%
}
\newcommand{\KalaTarjeta}[1]{%
  \begin{minipage}[t]{.235\linewidth}
  \setlength{\fboxsep}{3mm}%
  \fcolorbox{KalaLinea}{KalaPapel}{%
    \parbox[c][33mm][c]{\dimexpr\linewidth-2\fboxsep-2\fboxrule\relax}{%
      {\scriptsize\sffamily\color{KalaSuave}#1\hfill\csname KalaFamilia#1\endcsname}\par
      \centering\vspace{1mm}\KalaSimbolo[16mm]{#1}\par\vspace{1mm}
      {\small\sffamily\csname KalaNombre#1\endcsname}%
    }}%
  \end{minipage}%
}
\newcommand{\KalaCatalogoFila}[1]{%
  \KalaSimbolo[10mm]{#1}\quad\texttt{#1}&
  \textbf{\csname KalaNombre#1\endcsname}\par
  \csname KalaDescripcion#1\endcsname\par
  {\footnotesize\color{KalaSuave}\csname KalaFamilia#1\endcsname\;·\;
   \csname KalaTrayectos#1\endcsname\
   \ifnum\csname KalaTrayectos#1\endcsname=1 trayecto propuesto\else trayectos propuestos\fi.}\\[5pt]
}
% Inventario base de KalaScript, con las revisiones gráficas aceptadas.
% Coordenadas: cuadrícula 64 × 64, origen superior izquierdo.
% Editar directamente en LaTeX; no depende de SVG ni JSON.

\KalaRegistrar{K01}{Entidad}{Categoría}
  {Un objeto del dominio del que se habla; no identifica por sí solo un objeto particular.}{2}{%
  \draw (18,18) -- (18,40) .. controls (18,46) and (21.333333,49) .. (28,49) -- (44,49);
  \draw (30,15) -- (46,15) -- (46,31);
}

\KalaRegistrar{K02}{Propiedad}{Categoría}
  {Una característica atribuible a una entidad dentro de un dominio definido.}{2}{%
  \draw (16,44) .. controls (16,28) and (22,20) .. (34,20) .. controls (42,20) and (46,24.666667) .. (46,34) -- (32,34);
  \draw (32,34) -- (32,49);
}

\KalaRegistrar{K03}{Relación}{Categoría}
  {Con un índice y argumentos, aplica una relación del catálogo. El orden de los argumentos es significativo.}{2}{%
  \draw (13,42) .. controls (19,42) and (23.333333,38.333333) .. (26,31) .. controls (28.666667,23.666667) and (33,20) .. (39,20) -- (51,20);
  \draw (42,12) -- (51,20) -- (43,29);
}

\KalaRegistrar{K04}{Proceso}{Categoría}
  {Una actividad o transformación considerada en un dominio. No implica automáticamente causalidad.}{2}{%
  \draw (14,43) .. controls (14,26) and (30,26) .. (30,43) .. controls (30,58) and (47,49) .. (47,25);
  \draw (40,30) -- (47,22) -- (54,29);
}

\KalaRegistrar{K05}{Cantidad}{Categoría}
  {Una magnitud expresable mediante un valor y, cuando corresponde, una unidad.}{3}{%
  \draw (19,15) -- (19,48) -- (47,48);
  \draw (19,26) -- (36,26);
  \draw (19,37) -- (43,37);
}

\KalaRegistrar{K06}{Colección}{Categoría}
  {Una agrupación cuya naturaleza se declara: conjunto, lista u otra estructura. No impone orden ni unicidad por sí sola.}{2}{%
  \draw (24,14) .. controls (8,14) and (8,50) .. (24,50);
  \draw (34,22) -- (47,22) -- (34,42) -- (47,42);
}

\KalaRegistrar{K07}{Tiempo}{Categoría}
  {Un aspecto temporal, como instante o intervalo, con referencia temporal declarada cuando sea necesaria.}{2}{%
  \draw (17,16) .. controls (38.333333,16) and (49,21.333333) .. (49,32) .. controls (49,42.666667) and (38.333333,48) .. (17,48);
  \draw (18,25) -- (18,39) -- (29,39);
}

\KalaRegistrar{K08}{Espacio}{Categoría}
  {Un aspecto espacial cuya interpretación requiere dominio, escala o marco de referencia cuando corresponda.}{2}{%
  \draw (15,36) -- (32,15) -- (49,36) -- (32,49) -- cycle;
  % La barra termina en el centro del lado oblicuo, sin sobresalir.
  \draw (32,32) -- (45.761905,32);
}

\KalaRegistrar{K09}{Definición}{Marcador}
  {Introduce una definición nombrada. No afirma haber demostrado un resultado.}{3}{%
  \draw (14,18) -- (25,18) -- (25,46) -- (14,46);
  \draw (34,25) -- (51,25);
  \draw (34,39) -- (51,39);
}

\KalaRegistrar{K10}{Referencia}{Marcador}
  {Seguida de un índice, remite al referente definido en la leyenda del documento.}{1}{%
  \draw (17,44) -- (17,20) .. controls (17,16) and (19,14) .. (23,14) -- (41,14) -- (41,35) -- (29,35);
}

\KalaRegistrar{K11}{Negación}{Operador}
  {Niega la proposición de su grupo. No forma antónimos ni indica desconocimiento.}{2}{%
  \draw (15,19) .. controls (31,11) and (49,21) .. (49,36) .. controls (49,44) and (44.666667,48.333333) .. (36,49);
  \draw (19,50) -- (43,13);
}

\KalaRegistrar{K12}{Conjunción}{Operador}
  {Afirma conjuntamente dos proposiciones, sin afirmar un orden temporal entre ellas.}{2}{%
  \draw (14,20) -- (32,45) -- (50,20);
  \draw (32,45) -- (32,53);
}

\KalaRegistrar{K13}{Condicional}{Operador}
  {Implicación material clásica entre antecedente y consecuente. No significa causa ni ejecución.}{2}{%
  \draw (14,19) -- (32,32) -- (14,45);
  \draw (32,32) -- (50,32) -- (44,24);
}

\KalaRegistrar{K14}{Universal}{Operador}
  {Para cada elemento del dominio indicado se cumple el cuerpo agrupado.}{3}{%
  \draw (15,18) -- (15,45) .. controls (26.333333,51.666667) and (37.666667,51.666667) .. (49,45) -- (49,18);
  \draw (15,28) -- (49,28);
  \draw (32,12) -- (32,41);
}

\KalaRegistrar{K15}{Existencia}{Operador}
  {Existe al menos un elemento del dominio indicado para el que se cumple el cuerpo agrupado.}{2}{%
  \draw (15,18) -- (15,45) .. controls (26.333333,51.666667) and (37.666667,51.666667) .. (49,45) -- (49,18);
  \draw (25,28) -- (39,28);
}

\KalaRegistrar{K16}{Fuente}{Marcador}
  {Atribuye una proposición a una fuente identificada. No certifica su verdad ni cuantifica confianza.}{2}{%
  \draw (16,48) -- (16,16) -- (38,16) -- (49,27) -- (49,48) -- cycle;
  \draw (25,32) -- (39,32);
}

% Seis signos funcionales nuevos. Propuestas gráficas pendientes de prueba manual.
\KalaRegistrar{K17}{Alternativa}{Operador}
  {Disyunción inclusiva de dos proposiciones: al menos una es verdadera; pueden serlo ambas.}{2}{%
  \draw (14,18) -- (30,32) -- (14,46);
  \draw (50,18) -- (34,32) -- (50,46);
}
\KalaRegistrar{K18}{Posibilidad}{Operador}
  {Con un índice de modelo, afirma que el cuerpo se cumple en algún estado accesible. Fragmento modal proposicional.}{2}{%
  \draw (44,14) -- (20,14) .. controls (12,14) and (12,18) .. (12,26) -- (12,42) .. controls (12,50) and (16,50) .. (24,50) -- (44,50);
  \draw (32,23) -- (43,32) -- (32,41);
}
\KalaRegistrar{K19}{Pregunta}{Marcador}
  {Pregunta si la proposición agrupada es verdadera. Solo aparece en el nivel superior; no afirma el contenido.}{1}{%
  \draw (14,43) .. controls (14,16) and (45,14) .. (49,30) -- (33,40) -- (49,50);
}
\KalaRegistrar{K20}{Selección}{Constructor}
  {Construye el conjunto de elementos de un dominio que cumplen una condición. Vincula una variable en el cuerpo.}{2}{%
  \draw (14,16) -- (50,16) -- (35,34) -- (35,49);
  \draw (22,25) -- (42,25);
}
\KalaRegistrar{K21}{Secuencia}{Programa}
  {Ejecuta el primer programa y después el segundo. El segundo comienza solo si el primero termina.}{2}{%
  \draw (14,45) -- (14,32) -- (31,32) -- (31,19) -- (49,19);
  \draw (41,11) -- (49,19) -- (41,27);
}
\KalaRegistrar{K22}{Mientras}{Programa}
  {Comprueba una condición antes de cada iteración y ejecuta el cuerpo mientras sea verdadera. Puede no terminar.}{1}{%
  % Arco y barra en un solo trayecto; se elimina el segmento superpuesto.
  \draw (46,20) .. controls (20,5) and (9,33) .. (22,46) .. controls (34,58) and (54,41) .. (47,30) -- (37,30);
}

% Alias conservados para los comandos del cuaderno de comparación.
\KalaRegistrar{V06}{Colección B}{Aceptada}
  {Forma simplificada aceptada de K06. Alias gráfico de la forma habitual.}{2}{%
  \csname KalaDibujoK06\endcsname
}
\KalaRegistrar{V16}{Fuente B}{Aceptada}
  {Forma simplificada aceptada de K16. Alias gráfico de la forma habitual.}{2}{%
  \csname KalaDibujoK16\endcsname
}

% Formas anteriores, conservadas para leer y comparar documentos históricos.

\KalaRegistrar{H06}{Colección A}{Anterior}
  {Una agrupación cuya naturaleza se declara: conjunto, lista u otra estructura. No impone orden ni unicidad por sí sola.}{4}{%
  \draw (21,12) .. controls (10,12) and (13,24) .. (20,27) -- (14,32) -- (20,37) .. controls (13,40) and (10,52) .. (21,52);
  \draw (31,22) -- (46,22);
  \draw (31,32) -- (46,32);
  \draw (31,42) -- (46,42);
}

\KalaRegistrar{H16}{Fuente A}{Anterior}
  {Atribuye una proposición a una fuente identificada. No certifica su verdad ni cuantifica confianza.}{3}{%
  \draw (14,46) -- (14,19) -- (41,19) -- (41,46) -- cycle;
  \draw (22,11) -- (50,11) -- (50,37);
  \draw (22,29) -- (33,29);
}

% Construcciones escritas 0.2.0; las macros 0.1 siguen disponibles.
\newcommand{\KalaConcepto}[1]{\KalaSimbolo{K01}\kern-.8mm\textsubscript{#1}}
\newcommand{\KalaPropiedad}[1]{\KalaSimbolo{K02}\kern-.8mm\textsubscript{#1}}
\newcommand{\KalaVariable}[1]{\ensuremath{#1}}
\newcommand{\KalaColeccion}[2]{\KalaSimbolo{K06}\kern-.8mm\textsubscript{#1}\kern1mm\KalaGrupo{#2}}
\newcommand{\KalaTodo}[3]{\KalaSimbolo{K14}\kern1mm\ensuremath{#1\in}\,#2\kern1mm\KalaGrupo{#3}}
\newcommand{\KalaExiste}[3]{\KalaSimbolo{K15}\kern1mm\ensuremath{#1\in}\,#2\kern1mm\KalaGrupo{#3}}
\newcommand{\KalaAlternativa}[2]{\KalaSimbolo{K17}\kern1mm\KalaGrupo{#1\KalaSeparador#2}}
\newcommand{\KalaPosible}[2]{\KalaSimbolo{K18}\kern-.8mm\textsubscript{#1}\kern1mm\KalaGrupo{#2}}
\newcommand{\KalaPregunta}[1]{\KalaSimbolo{K19}\kern1mm\KalaGrupo{#1}}
\newcommand{\KalaSeleccion}[3]{\KalaSimbolo{K20}\kern1mm\ensuremath{#1\in}\,#2\kern1mm\KalaGrupo{#3}}
\newcommand{\KalaSecuencia}[2]{\KalaSimbolo{K21}\kern1mm\KalaGrupo{#1\KalaSeparador#2}}
\newcommand{\KalaMientras}[2]{\KalaSimbolo{K22}\kern1mm\KalaGrupo{#1\KalaSeparador#2}}
\newcommand{\KalaAsignar}[2]{\KalaSimbolo{K04}\kern-.8mm\textsubscript{1}\kern1mm\KalaGrupo{\KalaVariable{#1}\KalaSeparador\KalaMat{#2}}}
\newcommand{\KalaPosee}[2]{\KalaRelacion{8}{#1\KalaSeparador\KalaPropiedad{#2}}}
\newcommand{\KalaMiembro}[2]{\KalaRelacion{5}{#1\KalaSeparador#2}}
\newcommand{\KalaGusta}[2]{\KalaRelacion{7}{#1\KalaSeparador#2}}
\newcommand{\KalaFuenteB}[2]{\KalaSimbolo{V16}\kern1mm#1\kern1mm\KalaGrupo{#2}}
\newcommand{\KalaColeccionB}[2]{\KalaSimbolo{V06}\kern-.8mm\textsubscript{#1}\kern1mm\KalaGrupo{#2}}
\newcommand{\KalaFuenteAnterior}[2]{\KalaSimbolo{H16}\kern1mm#1\kern1mm\KalaGrupo{#2}}
\newcommand{\KalaColeccionAnterior}[2]{\KalaSimbolo{H06}\kern-.8mm\textsubscript{#1}\kern1mm\KalaGrupo{#2}}

% Banco único: identificador, área, intención, expresión, contexto, lectura, contraste, cobertura.
\newcommand{\KalaCaso}[8]{%
  \expandafter\def\csname KalaArea#1\endcsname{#2}%
  \expandafter\def\csname KalaIntencion#1\endcsname{#3}%
  \expandafter\def\csname KalaExpresion#1\endcsname{#4}%
  \expandafter\def\csname KalaContexto#1\endcsname{#5}%
  \expandafter\def\csname KalaLectura#1\endcsname{#6}%
  \expandafter\def\csname KalaContraste#1\endcsname{#7}%
  \expandafter\def\csname KalaCobertura#1\endcsname{#8}%
}
\newcommand{\KalaFicha}[1]{%
  \subsection*{#1 · \csname KalaArea#1\endcsname}
  \textbf{Intención.} \csname KalaIntencion#1\endcsname\par
  \KalaMostrar{\csname KalaExpresion#1\endcsname}
  \textbf{Contexto y leyenda.} \csname KalaContexto#1\endcsname\par
  \textbf{Lectura esperada.} \csname KalaLectura#1\endcsname\par
  \textbf{Contraste.} \csname KalaContraste#1\endcsname\par
  {\small\color{KalaSuave}\textbf{Cobertura.} \csname KalaCobertura#1\endcsname}\par
}

% Banco canónico de veinte pruebas. Las lecturas son soluciones del autor,
% no resultados de una prueba independiente de lectura.
\KalaCaso{T01}{Matemáticas}
 {Para cada natural hay un natural mayor.}
 {\KalaUniversal{x}{\mathbb{N}}{\KalaExistencia{y}{\mathbb{N}}{\KalaMat{x<y}}}}
 {\(\mathbb N=\{0,1,2,\ldots\}\). Las variables quedan vinculadas por sus cuantificadores.}
 {Para todo \(x\) natural existe \(y\) natural con \(x<y\); \(y\) puede depender de \(x\).}
 {Mover el existencial delante del universal exigiría un natural mayor que todos. La expresión actual es verdadera: sirve \(y=x+1\).}
 {control de compatibilidad con E01.}

\KalaCaso{T02}{Lógica y filosofía}
 {Negar que dos proposiciones sean ambas verdaderas.}
 {\KalaNegacion{\KalaConjuncion{\KalaMat{P}}{\KalaMat{Q}}}}
 {\(P,Q\) son proposiciones arbitrarias.}
 {No es verdad que \(P\) y \(Q\) sean ambas verdaderas.}
 {Con \(P=\text{verdadero}\) y \(Q=\text{falso}\), T02 es verdadera y T03 es falsa.}
 {negación con alcance sobre una combinación.}

\KalaCaso{T03}{Lógica y filosofía}
 {Afirmar que P es falsa y Q verdadera.}
 {\KalaConjuncion{\KalaNegacion{\KalaMat{P}}}{\KalaMat{Q}}}
 {Los mismos significados de \(P,Q\) que en T02.}
 {Se cumple la negación de \(P\) y se cumple \(Q\).}
 {El grupo de negación contiene solo \(P\). T02 y T03 no son equivalentes.}
 {comparación de alcances.}

\KalaCaso{T04}{Lectura y vida cotidiana}
 {A cada persona le gusta al menos un libro.}
 {\KalaTodo{x}{\KalaConcepto{1}}{\KalaExiste{y}{\KalaConcepto{2}}{\KalaGusta{\KalaVariable{x}}{\KalaVariable{y}}}}}
 {C1: personas del estudio; C2: libros del estudio. R7(persona, libro): a la persona le gusta ese libro según el criterio declarado.}
 {Cada persona puede elegir un libro diferente. No se exige un favorito único.}
 {Dos personas y dos libros, con gustos separados, hacen verdadera T04 y falsa T05.}
 {Extensión: conceptos indexados, dominios no numéricos y R7; cuantificadores conservados.}

\KalaCaso{T05}{Lectura y vida cotidiana}
 {Hay un mismo libro que les gusta a todas las personas.}
 {\KalaExiste{y}{\KalaConcepto{2}}{\KalaTodo{x}{\KalaConcepto{1}}{\KalaGusta{\KalaVariable{x}}{\KalaVariable{y}}}}}
 {Los mismos C1, C2 y R7 que en T04. En la comparación se usan dos personas y dos libros.}
 {Existe un libro del estudio que satisface la relación de gusto con cada persona.}
 {El testigo \(y\) se elige antes de recorrer las personas. No puede cambiar con \(x\).}
 {Extensión léxica; prueba de orden de cuantificadores.}

\KalaCaso{T06}{Arquitectura de computadoras}
 {El procesador lee el dato D de la dirección A.}
 {\KalaRelacion{2}{\KalaReferencia{1}\KalaSeparador\KalaReferencia{2}\KalaSeparador\KalaReferencia{3}}}
 {Referentes: 1 procesador; 2 dirección A; 3 dato D. Un único evento abstracto de lectura.}
 {El agente 1 lee el dato 3 desde la dirección 2.}
 {Intercambiar los argumentos 2 y 3 cambia sus papeles. No se describe una escritura ni un modelo de caché.}
 {relación ternaria R2.}

\KalaCaso{T07}{Biología}
 {La estructura A es parte propia de la célula B.}
 {\KalaRelacion{3}{\KalaReferencia{1}\KalaSeparador\KalaReferencia{2}}}
 {Referentes: 1 estructura A; 2 célula B, objetos concretos de un modelo ficticio. R3 es parte propia.}
 {A forma parte de B y A no es idéntica a B.}
 {No afirma que A sea un elemento de un conjunto B; esa es la relación de T08.}
 {distinción entre parte y miembro.}

\KalaCaso{T08}{Colecciones y matemáticas discretas}
 {El objeto A pertenece al conjunto formado por A y B.}
 {\KalaMiembro{\KalaReferencia{1}}{\KalaColeccion{1}{\KalaReferencia{1}\KalaSeparador\KalaReferencia{2}}}}
 {Referentes: 1 objeto A; 2 objeto B, distintos en este ejemplo. El índice 1 de K06 indica conjunto; el índice 2 indicaría lista.}
 {A es un miembro de \(\{A,B\}\).}
 {El conjunto \(\{A,B\}\) coincide con \(\{B,A\}\). Las listas \([A,B]\) y \([B,A]\) son diferentes. R5 solo acepta conjuntos o extensiones de conceptos.}
 {Extensión: construcción de colecciones y pertenencia R5.}

\KalaCaso{T09}{Historia}
 {El evento A terminó antes de comenzar B.}
 {\KalaRelacion{1}{\KalaReferencia{1}\KalaSeparador\KalaReferencia{2}}}
 {Referentes: 1 evento A, intervalo \([1,2]\); 2 evento B, intervalo \([3,4]\), misma escala temporal ficticia.}
 {El final de A precede estrictamente al inicio de B.}
 {La cronología no basta para obtener una afirmación causal como T10.}
 {relación temporal R1.}

\KalaCaso{T10}{Modelos causales}
 {En el modelo M, intervenir sobre X puede cambiar Y.}
 {\KalaRelacion{6}{\KalaReferencia{1}\KalaSeparador\KalaReferencia{2}\KalaSeparador\KalaReferencia{3}}}
 {Referentes: 1 variable X; 2 variable Y; 3 modelo determinista acíclico M con \(X:=U\), \(Y:=X\), \(U\in\{0,1\}\). Se compara el mismo contexto U bajo \(\operatorname{do}(X=0)\) y \(\operatorname{do}(X=1)\).}
 {En ese modelo, existe un contexto y dos intervenciones sobre X que producen valores distintos de Y.}
 {No afirma que una correlación observada demuestre causalidad, ni formula una teoría general de causa de eventos.}
 {Extensión: R6, con modelo y significado de intervención explícitos.}

\KalaCaso{T11}{Física experimental}
 {Atribuir una medición y su incertidumbre estándar a un registro.}
 {\KalaFuente{\KalaReferencia{2}}{\KalaMat{m(S_1)=5{,}0\,\mathrm{g};\quad u(m)=0{,}1\,\mathrm{g}}}}
 {Referente 2: registro ficticio de la muestra \(S_1\), con procedimiento y condiciones identificados. \(u(m)\) denota incertidumbre estándar.}
 {El registro 2 informa esos valores.}
 {K16 no certifica su verdad. La incertidumbre indicada no es un error conocido ni un intervalo de confianza.}
 {atribución y notación científica convencional.}

\KalaCaso{T12}{Alternativas lógicas}
 {Afirmar que P o Q se cumple, permitiendo ambas.}
 {\KalaAlternativa{\KalaMat{P}}{\KalaMat{Q}}}
 {\(P,Q\) son proposiciones; «o» se usa aquí en sentido inclusivo.}
 {Al menos una es verdadera.}
 {Si ambas son verdaderas, T12 sigue siendo verdadera. La exclusión se escribe combinando T12 con la negación de su conjunción.}
 {Extensión: K17, disyunción inclusiva.}

\KalaCaso{T13}{Posibilidad y modelos}
 {P es posible según el modelo modal 1.}
 {\KalaPosible{1}{\KalaMat{P}}}
 {Modelo 1: estados \(w_0,w_1\); estado actual \(w_0\); accesibilidad \(A_1=\{(w_0,w_1)\}\). P es falsa en \(w_0\) y verdadera en \(w_1\).}
 {Hay un estado accesible desde el actual donde P es verdadera.}
 {T13 es verdadera aunque P sea falsa en el estado actual. No asigna una probabilidad.}
 {Extensión: K18. Se limita al fragmento proposicional con modelo declarado.}

\KalaCaso{T14}{Preguntas y diálogo escrito}
 {Preguntar si el evento A precedió al B.}
 {\KalaPregunta{\KalaRelacion{1}{\KalaReferencia{1}\KalaSeparador\KalaReferencia{2}}}}
 {Referentes: 1 evento A; 2 evento B, con calendario común. En este ejercicio no se conocen sus fechas.}
 {¿Terminó A antes de comenzar B?}
 {Se solicita una respuesta; no se afirma la relación. «No sé» es un estado de información, distinto de negar R1.}
 {Extensión: K19. Pregunta polar en el nivel superior.}

\KalaCaso{T15}{Diseño de algoritmos}
 {Inicializar x en cero y aumentarlo mientras sea menor que tres.}
 {\KalaSecuencia{\KalaAsignar{x}{0}}{\KalaMientras{\KalaMat{x<3}}{\KalaAsignar{x}{x+1}}}}
 {Programa determinista sobre enteros matemáticos sin desbordamiento. La condición se comprueba antes del cuerpo.}
 {La traza de x después de inicializar es \(0,1,2,3\); después se detiene.}
 {Un bucle mientras con condición falsa ejecuta cero iteraciones. Una repetición hasta comprueba después de una primera ejecución. K13, el condicional lógico, no ejecuta pasos.}
 {Extensión: asignación indexada K04, secuencia K21 y mientras K22.}

\KalaCaso{T16}{Ingeniería y composición de conceptos}
 {Construir el concepto «sensores rojos y activos».}
 {\KalaSeleccion{x}{\KalaConcepto{3}}{\KalaConjuncion{\KalaPosee{\KalaVariable{x}}{1}}{\KalaPosee{\KalaVariable{x}}{2}}}}
 {C3: sensores del estudio. Propiedad 1: rojo según clasificación declarada. Propiedad 2: activo en el instante de inspección \(t_0\), común a toda la expresión.}
 {El conjunto de sensores que satisfacen ambas propiedades en el contexto indicado.}
 {Construye un concepto; no afirma que existan sensores que lo satisfagan. Para afirmar existencia se cuantifica sobre este resultado.}
 {Extensión: K20, propiedades indexadas y R8. Se reutilizan los conectivos.}

\KalaCaso{T17}{Química y notación especializada}
 {Registrar la proporción estequiométrica de formación de agua a partir de hidrógeno y oxígeno.}
 {\KalaMat{2\,\mathrm{H}_2+\mathrm{O}_2\longrightarrow 2\,\mathrm{H}_2\mathrm{O}}}
 {Región convencional de química. Los coeficientes indican proporciones estequiométricas; se omiten estados físicos y condiciones.}
 {Dos unidades estequiométricas de hidrógeno molecular y una de oxígeno molecular corresponden a dos de agua.}
 {La ecuación conserva cuatro átomos de H y dos de O a cada lado. No describe velocidad, mecanismo ni condiciones de realización.}
 {Interoperabilidad mediante región convencional; no se crea otro alfabeto químico.}

\KalaCaso{T18}{Geografía}
 {A está al norte de B en el marco geográfico F.}
 {\KalaRelacion{9}{\KalaReferencia{1}\KalaSeparador\KalaReferencia{2}\KalaSeparador\KalaReferencia{3}}}
 {Referentes: 1 lugar A con latitud \(12^\circ\); 2 lugar B con latitud \(10^\circ\); 3 marco F, mismo datum y época, latitud positiva hacia el norte. Lugares ficticios.}
 {La latitud de A es mayor que la de B en F.}
 {No afirma una distancia ni compara longitudes. Cambiar de marco exige declarar otro referente.}
 {Extensión: R9 con marco explícito.}

\KalaCaso{T19}{Economía descriptiva}
 {Atribuir a una tabla un aumento nominal del 20 por ciento.}
 {\KalaFuente{\KalaReferencia{1}}{\KalaMat{p_0=10\,\mathrm{UM};\quad p_1=12\,\mathrm{UM};\quad (p_1-p_0)/p_0=0{,}20}}}
 {Referente 1: tabla ficticia del precio del mismo bien, misma unidad de cantidad, en dos periodos. UM es una misma unidad monetaria; no hay ajuste por inflación.}
 {Según la tabla, el precio nominal pasa de 10 a 12 UM, un aumento relativo de 0,20.}
 {Un cambio nominal no establece un cambio real ni su causa. Es un ejemplo de representación de datos ficticios.}
 {fuente, cantidades y comparación convencional.}

\KalaCaso{T20}{Música y acústica}
 {Construir el concepto «tonos con frecuencia fundamental mayor que 440 Hz».}
 {\KalaSeleccion{x}{\KalaConcepto{5}}{\KalaMat{f(x)>440\,\mathrm{Hz}}}}
 {C5: tonos estables del estudio. \(f\) asigna a cada tono su frecuencia fundamental durante el mismo intervalo de observación.}
 {El conjunto de tonos del dominio cuya frecuencia fundamental supera 440 Hz.}
 {El umbral es estricto; 440 Hz queda fuera. No se selecciona por volumen. Un dominio vacío produce un conjunto vacío.}
 {Extensión: K20 con predicado convencional; reutilización de la misma composición en otro campo.}

